\section{核心问题：退化问题剖析}

为了系统性地研究深度增加对网络性能的影响，He等人设计了对照实验：在CIFAR-10数据集上分别训练一个20层的普通卷积网络（Plain Network）和一个56层的普通网络，两者均使用相同的卷积核尺寸和训练超参数。结果出人意料：56层网络在训练集上的误差（约12\%）比20层网络（约8\%）更高，测试集上的误差同样劣于20层网络。这种随深度增加性能反而下降的现象被定义为\textbf{退化问题}。值得注意的是，退化问题并非过拟合，因为过拟合通常表现为训练误差低但测试误差高，而这里训练误差本身也更高。因此，退化问题的根源在于优化困难：深层网络难以通过随机梯度下降法拟合恒等映射。

为什么深层网络难以学习恒等映射？理论上，如果浅层网络已经能够达到较好的性能，那么在其后添加恒等映射层（即输出等于输入）应至少不会使性能下降。然而，传统的堆叠层（Stacked Layers）中，网络需要学习一个映射 $H(x) = x$。当最优解确实接近恒等映射时，多层非线性变换的复合很难逼近单位变换。反向传播时，梯度需要穿越多个非线性激活函数（如ReLU），导致信号衰减或爆炸。退化问题表明，在现有优化框架下，让深层网络直接学习恒等映射比学习残差更加困难。

因此，解决退化问题的关键就是设计一种机制，使得网络能够更容易地学习恒等映射。残差学习正是针对这一需求提出的。它改变了目标函数的形式：不再期望底层堆叠层直接拟合 $H(x)$，而是拟合残差函数 $F(x) = H(x) - x$，从而将原始映射改写为 $H(x) = F(x) + x$。当 $H(x)$ 是恒等映射时，残差 $F(x)$ 只需趋于零，这比学习恒等映射本身简单得多。He等人的工作表明，通过引入快捷连接（Shortcut Connection）实现恒等映射的叠加，退化问题可以被有效消除，并且深度可以安全地扩展到数百层。